% F14 -- The whole ladder: gap to each run's own baseline, against training-set size.
\begin{figure}[t]
\centering
\plotlink{\nbfinal}{\fitwidth{%
\begin{tikzpicture}
\begin{axis}[rep, width=0.97\linewidth, height=6.4cm,
  xmode=log, log basis x=10,
  xmin=380, xmax=30000,
  xtick={400,900,1400,1900,5000,10000,15000,21000},
  xticklabels={400,900,1{,}400,1{,}900,5{,}000,10{,}000,15{,}000,21{,}000},
  x tick label style={font=\scriptsize,rotate=40,anchor=east},
  ymin=-0.5, ymax=7.6, ytick={0,1,...,5},
  yticklabels={0,$+1$,$+2$,$+3$,$+4$,$+5$},
  xlabel={training users (logarithmic scale)},
  ylabel={points gained over that run's own baseline},
  xmajorgrids, ymajorgrids,
]
% regime bands
\fill[clnonsig!30] (axis cs:380,-0.5) rectangle (axis cs:452,7.6);
\fill[clgain!12]   (axis cs:452,-0.5) rectangle (axis cs:1900,7.6);
\fill[clbase!12]   (axis cs:1900,-0.5) rectangle (axis cs:30000,7.6);
\node[font=\scriptsize,text=clrule,anchor=north,align=center] at (axis cs:900,7.45)
  {\textbf{GENERALISATION}\\ the four profiles become\\ large enough to read\\ \textcolor{clgain}{$+3.6\pt$}};
\node[font=\scriptsize,text=clrule,anchor=north,align=center] at (axis cs:8000,7.45)
  {\textbf{SATURATION}\\ every new user is a\\ behavioural duplicate\\ \textcolor{clbase}{$+0.1\pt$}};
\node[font=\scriptsize,text=clrule,anchor=south,align=center,rotate=90] at (axis cs:414,1.2)
  {\textbf{IMITATION}};
%
\addplot[clmodel,thick,mark=*,mark size=1.8pt] coordinates {
(400,0.6)(900,2.8)(1400,3.5)(1900,4.2)(5000,4.3)(10000,4.4)(15000,4.7)(21000,4.3)};
\addplot[clbase,dashed,thick,domain=380:30000,samples=2]{0};
%
\node[font=\scriptsize,anchor=west] at (axis cs:432,0.95) {$+0.6$};
\node[font=\scriptsize,anchor=north] at (axis cs:900,2.62) {$+2.8$};
\node[font=\scriptsize,anchor=north] at (axis cs:1400,3.32) {$+3.5$};
\node[font=\scriptsize,anchor=north west] at (axis cs:1950,4.15) {$+4.2$};
\node[font=\scriptsize,anchor=south] at (axis cs:5000,4.35) {$+4.3$};
\node[font=\scriptsize,anchor=south] at (axis cs:10000,4.45) {$+4.4$};
\node[font=\scriptsize,anchor=south] at (axis cs:15000,4.75) {$+4.7$};
\node[font=\scriptsize,anchor=south] at (axis cs:21000,4.35) {$+4.3$};
\node[font=\scriptsize,text=clbase,anchor=east] at (axis cs:28000,0.24) {the persistence baseline of each run};
\end{axis}
\end{tikzpicture}}}
\caption{The whole scaling ladder, in points over each run's own baseline}
\label{fig:scaling-ladder}
\figtext{%
Eight training-set sizes, each read against the persistence baseline measured
on its \emph{own} test users. The horizontal axis is logarithmic, so equal
distances are equal multiples and 400 to 1{,}900 occupies the same width as
4{,}000 to 19{,}000; the vertical axis is not accuracy but the gap to the
baseline, the only quantity that survives a change of test set, so the dashed
line at zero means ``indistinguishable from repeating the last cell''. At 400
training users the model gains $+0.6\pt$ and is barely distinguishable from the rule it
is meant to beat. Between 400 and 1{,}900 the edge climbs $+0.6 \rightarrow
+2.8 \rightarrow +3.5 \rightarrow +4.2\pt$, a four-fold increase in population
buying $+3.6\pt$, which is essentially the entire benefit the project ever
obtained from data volume. The ten-fold increase that follows buys $+0.1\pt$,
and not monotonically: the best tier (15{,}000) is followed by the worst
(21{,}000), and none of the steps inside that stretch is statistically
distinguishable from zero (Table~\ref{tab:bootstrap}). The curve does not slow
down, it stops. Two things constrain how far that reads. The 900 to
1{,}900-user tiers are scored on \stwok{} (100 test users, 48.2\,\% baseline)
and the 400-user point is the \smob{} run (43.1\,\% baseline) and the 5{,}000 to
21{,}000-user tiers on \seleven{} (1{,}000 test users, 50.2\,\% baseline),
which is why the axis is a gap and why the small step across that boundary is
not a measurement. And every tier received the same budget of optimiser steps,
so at 21{,}000 users only 36\,\% of the available training windows were ever
read: the defensible claim is that 1{,}900 users are already enough at this
training budget, not that more users can never help.}
\figsource{\nbfinal}{train\_cellid\_final.ipynb}{the eight-tier scaling notebook}
\end{figure}
